% !Mode:: "TeX:UTF-8"
\documentclass{mcmthesis}
%\documentclass[CTeX = true]{mcmthesis}  % 当使用 CTeX 套装时请注释上一行使用该行的设置
\setlength{\headheight}{15pt}
\mcmsetup{tstyle=\color{red}\bfseries,%修改题号，队号的颜色和加粗显示，黑色可以修改为 black
	tcn = 2604335, problem = B, %修改队号，参赛题号
	sheet = true, titleinsheet = true, keywordsinsheet = true,
	titlepage = false, abstract = true}
%四款字体可以选择
%\usepackage{times}
%\usepackage{newtxtext}
%\usepackage{palatino}
\usepackage{txfonts}
\usepackage{tabularx}
\usepackage{booktabs}
\usepackage{textcomp}
\usepackage{makecell}
\usepackage{booktabs}  % 如果你用了 \bottomrule
\usepackage{enumitem}
\usepackage{float}
\usepackage{cite}
\usepackage {caption}
\usepackage{indentfirst}  %首行缩进，注释掉，首行就不再缩进。
\usepackage{lipsum}
\title{LATEX TITLE}
\begin{document}
	\begin{abstract}
		\par Use this template to begin typing the first page (summary page) of your electronic report. This template uses a 12-point Times New Roman font. Submit your paper as an Adobe PDF electronic file (e.g. 1111111.pdf), typed in English, with a readable font of at least 12-point type.
		
		Do not include the name of your school, advisor, or team members on this or any page.
		
		Papers must be within the 25 page limit.
		
		Be sure to change the control number and problem choice above.
		You may delete these instructions as you begin to type your report here.
		
		Follow us @COMAPMath on Twitter or COMAPCHINAOFFICIAL on Weibo for the most up to date contest information.
		
		\begin{keywords}
			keyword1; keyword2
		\end{keywords}
	\end{abstract}
	\maketitle
	%% Generate the Table of Contents, if it's needed.
	\tableofcontents
	\newpage
	%%
	%% Generate the Memorandum, if it's needed.
	%% \memoto{\LaTeX{}studio}
	%% \memofrom{Liam Huang}
	%% \memosubject{Happy \TeX{}ing!}
	%% \memodate{\today}
	%% \memologo{\LARGE I'm pretending to be a LOGO!}
	%% \begin{memo}[Memorandum]
		%%   \lipsum[1-3]
		%% \end{memo}
	%%
	\section{Introduction}
	\subsection{Problem Background}
    \lipsum[3]
	\subsection{Restatement of the Problem}
	\lipsum[3]
	\subsection{Literature Review}
	\lipsum[3]
	\begin{itemize}
		\item the angular velocity of the bat,
		\item the velocity of the ball, and
		\item the position of impact along the bat.
	\end{itemize}
	\lipsum[4]
	\emph{center of percussion} [Brody 1986], \lipsum[5]
	
	\begin{Theorem} \label{thm:latex}
		\LaTeX
	\end{Theorem}
	\begin{Lemma} \label{thm:tex}
		\TeX .
	\end{Lemma}
	\begin{proof}
		The proof of theorem.
	\end{proof}
	
	\subsection{Our Work}
	\lipsum[6]
	\begin{itemize}
		\item
		\item
		\item
		\item
	\end{itemize}
	
	\lipsum[7]
	
	\section{Assumptions and Justifications}
	\begin{itemize}
	\item \textbf{Assumption 1}:
	\item\textbf{Justification}:
	\item \textbf{Assumption 2}:
	\item\textbf{Justification}
	\end{itemize}
	\section{Notations}
	The primary symbols used in our model are categorized into deterministic parameters and stochastic variables, as summarized in Tables~\ref{tab:deterministic} and \ref{tab:stochastic}.
	
	\begin{table}[htbp]
		\centering
		\caption{Deterministic Parameters and Decision Variables}
		\label{tab:deterministic}
		\begin{tabular}{l p{9cm} c} 
			\toprule
			\textbf{Symbol} & \textbf{Description} & \textbf{Unit} \\
			\midrule
			$M_{\text{total}}$      & Total material demand for Moon Colony construction ($10^8$ MT) & MT \\
			$x, y$                  & Payload mass allocated to Space Elevator System (SES) and Rocket System (RS), respectively & MT \\
			$R_e$                   & Total annual transport capacity of the three Galactic Harbours (SES) & MT/year \\
			$R_r$                   & Total annual transport capacity of the rocket system & MT/year \\
			$M_{L\text{-}ele}$      & Annual payload capacity per Galactic Harbour & MT/year \\
			$M_{R\text{-}rocket}$   & Payload capacity per rocket launch (2050 projection) & MT/launch \\
			$f_{\text{rocket}}$     & Annual launch frequency per rocket site & launches/year \\
			$C_{ele}$               & Unit transport cost via space elevator (to apex anchor) & \$/MT \\
			$C_{\text{rocket}}$     & Unit transport cost via rocket (Earth to Moon) & \$/MT \\
			$T$                     & Total delivery time for construction materials & year \\
			$C$                     & Total accumulated logistics cost & \$ \\
			$R_{\text{gear}}$       & Mass penalty factor for lunar descent propulsion (gear ratio) & -- \\
			$C_{\text{repair}}$     & Additional repair cost per elevator failure & \$ \\
			$C_{\text{maint/yr}}$   & Extra annual maintenance cost induced by failures & \$/year/harbour \\
			\bottomrule
		\end{tabular}
	\end{table}
	
	\begin{table}[htbp]
		\centering
		\caption{Key Stochastic Variables for Robustness Analysis}
		\label{tab:stochastic}
		\begin{tabular}{l p{9cm} c} 
			\toprule
			\textbf{Symbol} & \textbf{Description} & \textbf{Unit} \\
			\midrule
			$\tilde{\kappa}$            & Effective capacity factor due to tether swaying ($\tilde{\kappa} \sim \mathrm{Beta}(8,2)$, $\mathbb{E}[\tilde{\kappa}] = 0.80$) & -- \\
			$q$                         & Rocket failure probability per launch ($q \sim \mathrm{Beta}(2,98)$, $\mathbb{E}[q] = 0.02$) & -- \\
			$\tilde{T}_A, \tilde{T}_B, \tilde{T}_C$ & Stochastic delivery time for Scenarios A (pure elevator), B (pure rocket), and C (hybrid) & year \\
			$\tilde{C}_A, \tilde{C}_B, \tilde{C}_C$ & Stochastic total cost for Scenarios A, B, and C & \$ \\
			\bottomrule
		\end{tabular}
	\end{table}
	\begin{table}[htbp]
		\centering
		\caption{Key Variables for Water Logistics Analysis}
		\label{tab:water}
		\begin{tabular}{l p{9cm} c} 
			\toprule
			\textbf{Symbol} & \textbf{Description} & \textbf{Unit} \\
			\midrule
			$N$                     & Lunar colony population (100,000 persons) & person \\
			$Q_{\text{daily}}$      & Daily water consumption per capita ($\approx 4.53$ kg/person/day, NASA) & kg/person/day \\
			$\eta$                  & Water recycling efficiency ($0\% \sim 99\%$) & -- \\
			$M_{\text{net\_water}}$ & Annual net water import requirement after recycling & ton/year \\
			$x_{\text{water\_ele}}$ & Water mass allocated to space elevator transport & ton/year \\
			$x_{\text{water\_roc}}$ & Water mass allocated to rocket transport (overflow) & ton/year \\
			\bottomrule
		\end{tabular}
	\end{table}
	\section{The Preparation of Model}
	\subsection{Data Overview}
	In this subsection, we compile the key mission parameters supplied by the MCM, forming the basis for all subsequent computational analyses. These constants serve as the boundary conditions for the POHL model.
	
	\begin{itemize}[leftmargin=*]
		\item \textbf{Material Demand ($M_{total}$):} The Moon Colony requires a total of $10^8$ metric tons (MT) of materials for infrastructure and construction.
		\item \textbf{Individual SES Capacity ($M_{L-ele}$):} Each Galactic Harbour, featuring a dual-tether system, provides a baseline transport capacity of 179,000 MT annually.
		\item \textbf{Rocket Payload ($M_{R-rocket}$):} By 2050, advanced heavy-lift launch vehicles (e.g., evolved Falcon Heavy or Starship variants) are expected to deliver between 100 and 150 MT per flight directly to the lunar surface.
		\item \textbf{Target Population ($P_{pop}$):} The logistics model is designed to support a sustained population of $10^5$ residents starting in the year 2050.
	\end{itemize}
	\subsection{Data Collection}
		Since certain key data are still lacking in the research context, we have collected the required information from a variety of official and authoritative data sources, laying a solid foundation for the subsequent extrapolation of data to 2050.\\
\begin{table}[htbp]
	\centering
	\caption{Key Data for Self-Sustaining Moonbase}
	\label{tab:moonbase_data_source}
	% 两列表格，自适应页面宽度
	\begin{tabularx}{\linewidth}{lX}
		\toprule
		\textbf{Data} & \textbf{Description} \\
		\midrule
		2.3 acre-ft/year & Domestic water demand (100 people/year) \cite{gerald20xxwater} \\
		72 acre-ft/year & Agricultural water demand (100 people/year) \cite{gerald20xxwater} \\
		2 acre-ft/year & Technical water demand (cooling/fire protection) \cite{gerald20xxwater} \\
		98\% (\textpm5\%) & ISS water recycling efficiency (post-BPA) \cite{wil20xxiss} \\
		\$100/150/kg & Space elevator: pure operation cost (electricity + maintenance)\cite{isec2019mars} \\
		\$200/kg & Space elevator: total cost (operation + depreciation + risk premium)\cite{isec2019mars} \\
		\$405/kg & Rocket fuel cost only (elevator top to Moon, recyclable rocket) \cite{isec2019mars}\\
		\$505/kg & Total cost (space elevator + recyclable rocket fuel) \cite{isec2019mars}\\
		2.5–3.0 kg CO/kg & RP-1 combustion emission factor \cite{ipcc2014ar5} \\
		820–1,050 g CO-eq/kWh & Coal power carbon intensity (global average) \cite{ipcc2014ar5} \\
		11–12 g CO-eq/kWh & Onshore wind power carbon intensity \cite{ipcc2014ar5} \\
		17,860 tons CO/year & Annual CO emissions (global avg. grid, 37.6 GWh) \cite{ipcc2014ar5} \\
		451 tons CO/year & Annual CO emissions (100\% wind, 37.6 GWh) \cite{ipcc2014ar5} \\
		37,600 tons CO/year & Annual CO emissions (100\% coal, 37.6 GWh) \cite{ipcc2014ar5} \\
		\bottomrule
	\end{tabularx}
\end{table}
	\subsection{Data Extrapolation for the 2050 Scenario}
	\subsubsection{Forecast of Rocket Launch Frequency via ARIMA Model}We extrapolated the annual number of rocket launches for 2050 using the ARIMA(p, d, q) model, which is adopted to address the distinct time dependence and non-stationary trends in global rocket launch data. By removing long-term trends through differencing (d) and combining autoregressive (AR) and moving average (MA) components to capture historical fluctuations, this model better characterizes the growth inertia of the aerospace industry than traditional trend-fitting approaches.\\
	
	% TODO: \usepackage{graphicx} required
	\begin{figure}[H]
		\centering
		\includegraphics[width=0.7\linewidth]{../ARIMA}
		\caption{}
		\label{fig:arima}
	\end{figure}
	As illustrated in Figure 1, the model projects a consistent upward trend in global space activity. By the target year of 2050, the annual launch frequency is estimated to reach 2,253, serving as a pivotal baseline for calculating total lunar payload delivery and assessing the corresponding environmental footprint.\\
	\subsubsection{Cost Projection of Advanced Rockets via Wright’s Learning Law}
	\subsection{Cost Projection of Advanced Rockets via Wright's Learning Law}
	To assess the economic feasibility of Scenario B, we employ \textbf{Wright's Law} to project the unit launch cost in 2050. This model accounts for the significant cost reduction driven by the economies of scale from the $10^8$ metric ton lunar colony requirement.\\
	\subsubsection{LEO Unit Cost Derivation}
	The unit cost to Low Earth Orbit (LEO) is modeled by the learning curve equation:
	\begin{equation}
		C_{LEO, 2050} = C_1 \cdot X^{-b}
	\end{equation}
	Where:
	\begin{itemize}
		\item $C_1 = \$1,520/\text{kg}$ is the baseline cost of Falcon Heavy in 2023.
		\item $X$ represents the cumulative doubling factor of space lift capacity.
		\item $b$ is the learning index. Based on a standard aerospace learning rate of 88\%, we set $b \approx 0.184$.
	\end{itemize}
	By extrapolating the cumulative production growth to meet the 2050 mission scale, the model yields:
	\begin{equation}
		C_{LEO, 2050} \approx \$380/\text{kg}
	\end{equation}
	To account for the additional energy and hardware required for Trans-Lunar Injection (TLI) and powered descent, we introduce the Lunar Transfer Multiplier ($\lambda = 7.2$). The final projected cost to the lunar surface is:
	\begin{equation}
		C_{Moon} = C_{LEO, 2050} \times \lambda = 380 \times 7.2 = \$2,736/\text{kg}
	\end{equation}
	This cost is a key parameter for the subsequent comparative analysis between traditional rocket logistics and the space elevator system.
	\section{The POHL Model: Pareto-Optimal Hybrid Logistics Framework}
	% TODO: \usepackage{graphicx} required
	\begin{figure}[H]
		\centering
		\includegraphics[width=0.8\linewidth]{"../Figure 2 Operational Flowchart of Lunar Infrastructure Transport Strategy Optimization"}
		\caption{}
		\label{fig:figure-2-operational-flowchart-of-lunar-infrastructure-transport-strategy-optimization}
	\end{figure}
	\subsection{Mathematical Formulation of the Logistics System}
	The construction of a lunar colony is a multi-dimensional challenge characterized by the conflict between temporal efficiency and economic sustainability. The \textbf{Pareto-Optimal Hybrid Logistics (POHL) Model} is developed to navigate this trade-off by optimizing the resource allocation between the Space Elevator System (SES) and the Rocket System (RS).
	
	\subsubsection{Decision Variables and Mass Balance Constraints}
	We define $x$ and $y$ as the primary decision variables, representing the total mass allocated to the SES and RS, respectively. To ensure the complete delivery of the required infrastructure, the following mass balance constraint must be satisfied:
	\begin{equation}
		x + y = M_{total}
	\end{equation}
	where $M_{total} = 10^8$ metric tons (MT), with the non-negativity constraints $x, y \geq 0$.
	
	\subsubsection{Cost Modeling and the Physics-based "Gear Ratio" Correction}
	A realistic economic assessment of lunar logistics must account for the propulsive requirements of the final landing phase. We introduce the \textbf{Gear Ratio} ($R_{gear}$) to adjust the unit cost of the SES. While the elevator delivers cargo to the Apex Anchor with near-zero fuel consumption, the subsequent Lunar Orbit Insertion (LOI) and Powered Descent still necessitate chemical propulsion. 
	
	The comprehensive unit cost for the SES ($C_{ele}$) is formulated as:
	\begin{equation}
		C_{ele} = \underbrace{\left( \frac{\sum I_{C\,APEX}}{L \cdot \sum K_{annual}} + C_{ops} + C_{maint} \right)}_{\text{Levelized Cost of Ascent to Apex}} \times \underbrace{R_{gear}}_{\text{Mass Penalty}}
	\end{equation}
	where $R_{gear} \approx 2.2$ serves as the correction factor derived from the Tsiolkovsky rocket equation for lunar landing scenarios.
	
	\subsubsection{Formulation of Competing Objectives}
	The system performance is evaluated through two objective functions:
	\begin{itemize}
		\item \textbf{Temporal Efficiency ($T$):} Assuming a high degree of operational parallelism, the total mission duration is governed by the system's throughput bottleneck:
		\begin{equation}
			T(x) = \max \left( \frac{x}{R_e}, \frac{M_{total} - x}{R_r} \right)
		\end{equation}
		where $R_e$ and $R_r$ represent the total annual transport capacities of the SES and RS, respectively.
		\item \textbf{Economic Expenditure ($C$):} The total fiscal requirement is the linear summation of costs:
		\begin{equation}
			C(x) = C_{ele} \cdot x + C_{rocket} \cdot (M_{total} - x)
		\end{equation}
	\end{itemize}
	
	\subsection{Analytical Solution and Pareto Derivation}
	
	\subsubsection{Boundary Baseline Evaluation}
	Evaluating the boundary conditions provides a critical baseline for the optimization space:
	\begin{itemize}
		\item \textbf{Scenario A ($x=M_{total}$):} Utilizing only the SES yields a total duration of $T_A \approx 186.2$ years. This excessive timeline is strategically unviable for a colony intended to stabilize by 2050.
		\item \textbf{Scenario B ($x=0$):} Utilizing only the RS provides the minimum possible duration but at an astronomical cost, potentially exceeding budgetary limits.
	\end{itemize}
	
	\subsubsection{Derivation of the Analytical Pareto Frontier}
	By eliminating the decision variable $x$ between the objective functions, we derive the \textbf{Analytical Pareto Frontier} for Scenario C, defining the set of all non-dominated solutions:
	\begin{equation}
		C(T) = R_e(C_{ele} - C_{rocket})T + C_{rocket} M_{total}
	\end{equation}
	Since $C_{ele} < C_{rocket}$, the frontier is a line with a negative slope, indicating a constant \textbf{Marginal Rate of Substitution (MRS)} between budget and schedule.
	
	\subsubsection{Determination of the Optimal Strategy through Robust Compromise}
	To identify a robust strategy, we first solve for the time-optimal point ($x_{time}^*$), where both delivery modes conclude their tasks simultaneously:
	\begin{equation}
		x_{time}^* = M_{total} \cdot \frac{R_e}{R_e + R_r}
	\end{equation}
	Adopting the geometric midpoint of the rational decision space between $x_{time}^*$ and $x_{cost}^* = M_{total}$, we achieve the balanced hybrid strategy ($x_{mixed}$):
	\begin{equation}
		x_{mixed} = M_{total} \cdot \frac{2R_e + R_r}{2(R_e + R_r)}
	\end{equation}
	The optimized total cost is then expressed as a weighted function of the two transport modes:
	\begin{equation}
		C_{mixed} = \frac{M_{total}}{2(R_e + R_r)} \left[ C_{ele}(2R_e + R_r) + C_{rocket} R_r \right]
	\end{equation}
	\subsection{Solution and Result Analysis}
	Based on the defined parameters ($M_{total}=10^8$ MT, $R_e=537,000$ MT/unit, $R_r=281,625$ MT/unit), we evaluate the logistics performance across three boundary and optimized scenarios.
	
	\subsubsection{Boundary Baselines: Scenario A and B}
	\begin{itemize}
		\item \textbf{Scenario A (Pure SES)}: $x = M_{total}, y = 0$. This configuration minimizes costs to \textbf{\$60.50 Trillion} but extends the duration to \textbf{186.22 units}.
		\item \textbf{Scenario B (Pure RS)}: $x = 0, y = M_{total}$. Utilizing only the rocket fleet results in the highest expenditure of \textbf{\$273.60 Trillion} and a duration of \textbf{355.08 units}.
	\end{itemize}
	\hspace{1.5em} The simulation results reveal that $T_B$ (355 units) is significantly higher than $T_A$ (186 units), suggesting that the aggregate rocket throughput $R_r$ is relatively constrained compared to the elevator system. This disparity stems from the fact that rocket operations are strictly governed by optimal launch windows and complex propellant logistics. As a result, the total logistics flux of the Rocket System (RS) in 2050 remains lower than the throughput of a mature Space Elevator System (SES) infrastructure.\\
	\subsubsection{Optimized Benchmarks: Time-Optimal and POHL Mixed Strategy}
	To resolve the efficiency-cost conflict, we calculate the time-optimal pivot and the final POHL compromise:
	\begin{enumerate}
		\item \textbf{Time-Optimal Solution ($x_{time}^*$)}: By balancing throughput ($x/R_e = y/R_r$), we achieve the minimum mission duration of \textbf{122.16 units} with $x \approx 65.60$ MT and $y \approx 34.40$ MT, totaling \textbf{\$133.81 Trillion}.
		\item \textbf{POHL Mixed Strategy (Scenario C)}: Implementing the geometric compromise $x_{mixed} = M_{total} \cdot \frac{2R_e + R_r}{2(R_e + R_r)}$, we derive the following optimal allocation:
		\begin{itemize}
			\item \textbf{SES Allocation ($x_{mixed}$)}: $82,798,900.60$ MT (82.8\%).
			\item \textbf{RS Allocation ($y_{mixed}$)}: $17,201,099.40$ MT (17.2\%).
			\item \textbf{System Performance}: $T_{mixed} = \mathbf{154.19}$ \textbf{units}, $C_{mixed} = \mathbf{\$97.16}$ \textbf{Trillion}.
		\end{itemize}
	\end{enumerate}
	% TODO: \usepackage{graphicx} required
	\begin{figure}[H]
		\centering
		% 第一张图：Pareto Frontier
		\begin{minipage}[b]{0.6\linewidth}
			\centering
			% 使用 trim 手动切除白边：左 下 右 上 (根据预览微调数值)
			% clip 必须设为 true 裁剪才会生效
			\includegraphics[width=1\linewidth, trim=10pt 10pt 10pt 10pt, clip]{"../Pareto frontier"}
			\caption{Pareto Frontier of Logistics Strategy}
			\label{fig:pareto-frontier}
		\end{minipage}
		\hfill % 在两个 minipage 之间加入弹性间距，推向两端
		% 第二张图：Allocation (堆叠柱状图)
		\begin{minipage}[b]{0.35\linewidth}
			\centering
			% 同样的，如果 Allocation 也有白边，加入 trim 参数
			\includegraphics[width=1\linewidth, trim=5pt 5pt 5pt 5pt, clip]{../Allocation}
			\caption{Allocation Share}
			\label{fig:allocation}
		\end{minipage}
	\end{figure}
	
	The POHL Mixed Strategy (Scenario C) successfully mitigates the extreme delays of Scenario A and the exorbitant costs of Scenario B, providing a robust baseline for the subsequent stochastic analysis.\\
	\section{Resilience of the POHL Model: Robustness Analysis under Stochastic Disturbances}
In Problem 2, we transition from a deterministic framework to a stochastic paradigm to evaluate system performance under operational uncertainties, specifically focusing on \textbf{Tether Swaying}, \textbf{Rocket Failure}, and \textbf{Elevator Breakdown}.
\subsection{Stochastic Characterization of Disruption Factors}

\subsubsection{Coriolis-induced Tether Swaying}
To account for capacity reductions caused by safety protocols (e.g., speed reduction to dampen oscillations), we introduce the \textbf{Swaying Factor} $\tilde{\kappa} \sim Beta(8, 2)$. 
The expected efficiency $\mathbb{E}[\tilde{\kappa}] = 0.8$ reflects the engineering reality where the system typically operates at 80\% of its theoretical maximum capacity. This left-skewed distribution captures the conservative nature of operational throughput.

\subsubsection{Discrete Rocket Failure}
Rocket launches are modeled as independent Bernoulli trials. Under a Bayesian framework, the failure probability $p$ follows a prior distribution $p \sim Beta(2, 98)$, yielding an expected failure rate of 2\%. For a required $N_{req}$ successful missions, the actual number of launches $\tilde{N}_{act}$ follows a \textbf{Negative Binomial distribution}:
\begin{equation}
	\tilde{N}_{act} \sim \mathrm{NB}(N_{req}, \tilde{p})
\end{equation}

\subsubsection{Continuous Elevator Breakdown}
System failures are modeled as a \textbf{Compound Poisson Process}. The number of failures $N(t)$ within time $t$ follows $\mathrm{Poisson}(\lambda t)$, with $\text{MTBF} = 5 \times 10^4$ h. Each repair duration $\tilde{D}_j$ is governed by an \textbf{Exponential distribution} $\tilde{D}_j \sim \mathrm{Exp}(\mu)$ with $\text{MTTR} = 720$ h. The cumulative downtime is defined as $T_{down}(t) = \sum_{j=1}^{N(t)} \tilde{D}_j$.

\subsection{Model Reformulation under Stochasticity}

We modify the original transport models by incorporating these random variables into the time and cost functions.

\subsubsection{Scenario A: Pure SES Transport (SES-Only)}
The total duration $\tilde{T}_A$ is reformulated as an additive model of nominal operation and repair time:
\begin{equation}
	\tilde{T}_{A} = \frac{M_{total}}{R_e \cdot \tilde{\kappa}} + \sum_{j=1}^{\tilde{k}_{fail}} \tilde{D}_j, \quad \text{where } \tilde{k}_{fail} \sim \mathrm{Pois}(\lambda \cdot t_{nom})
\end{equation}
The total cost $\tilde{C}_A$ includes variable maintenance and discrete repair costs $C_{rep} \approx \$100M$:
\begin{equation}
	\tilde{C}_{A} = M_{total} \cdot C_{ele} + C_{rep} \cdot \tilde{k}_{fail} + C_{maint/yr} \cdot \tilde{T}_{A}
\end{equation}

\subsubsection{Scenario B: Pure Rocket Transport (RS-Only)}
Relying solely on the rocket fleet, the performance is driven by the realization of $\tilde{N}_{act}$:
\begin{equation}
	\tilde{T}_{B} = \frac{\tilde{N}_{act}}{10 \cdot f_{rocket}}, \quad \tilde{C}_{B} = \tilde{N}_{act} \cdot C_{rocket}
\end{equation}

\subsubsection{Scenario C: The Hybrid POHL Strategy}
The hybrid model leverages a \textbf{Diversification Effect}. The total duration is governed by the bottleneck of parallel sub-systems:
\begin{equation}
	\tilde{T}_{C} = \max \left( \tilde{T}_{ele}(x_{mixed}), \, \tilde{T}_{rocket}(y_{mixed}) \right)
\end{equation}
The stochastic cost $\tilde{C}_C$ aggregates the expenses from both modes, adjusted for random failures.

\subsection{Monte Carlo Simulation and Probabilistic Insights}

We implemented $10^5$ iterations to evaluate the statistical distributions of the three scenarios. 

\subsubsection{Robustness and Risk Hedging}
Numerical results indicate that Scenario C achieves a lower \textbf{Coefficient of Variation (CV)} compared to A or B. This stems from the fact that rocket and elevator risks are imperfectly correlated. Mathematically, the variance of the hybrid strategy $\sigma_{C}^2 \approx w^2 \sigma_{A}^2 + (1-w)^2 \sigma_{B}^2$ is suppressed, demonstrating a classic \textbf{Portfolio Hedging Effect}.

\subsubsection{Stochastic Dominance and Stability}
The simulation confirms that Scenario C maintains \textbf{Stochastic Dominance} with a probability $P > 75\%$ even under $\pm20\%$ parameter fluctuations. While Scenario A suffers from "thick-tail" time risks and Scenario B from exorbitant budget volatility, Scenario C remains the most resilient. In the presence of real-world failures, the mixed strategy transcends being a mere compromise; it functions as an essential hedge against high-impact, low-probability "Black Swan" events in space logistics.
\section{Sustainable Colony Metabolism: The POHL Framework under Vital Constraints}

Problem 3 shifts from "infrastructure delivery" to "steady-state survival." We define the colony as a biological entity with a specific \textbf{Metabolic Load}, evaluating whether the established POHL framework can sustain $10^5$ residents without triggering a logistics collapse.

\subsection{Metabolic Demand Function: The Water-Loop Sensitivity}
The survival of the colony is parameterized by the Environmental Control and Life Support System (ECLSS) efficiency. Based on NASA's biological requirements, the annual net water inflow required from Earth, $M_{net\_water}$, is a function of the recycling rate $\eta$:
\begin{equation}
	M_{net\_water}(\eta) = \underbrace{N \cdot Q_{daily} \cdot 365}_{\text{Gross Demand}} \cdot \underbrace{(1 - \eta)}_{\text{System Loss}} \cdot \underbrace{(1 + \alpha_{tank})}_{\text{Structure Penalty}}
\end{equation}
Using $N=10^5$ and $Q_{daily}=4.53$ kg/day, the demand is extremely sensitive to $\eta$. This establishes the \textbf{Metabolic Load} that must be injected into the existing delivery loops.

\subsection{Priority-based Resource Allocation and Overflow Logic}
We introduce a \textbf{Hierarchical Allocation Model} where limited transport capacity is treated as a finite container. Logistics are prioritized as follows:
\begin{enumerate}
	\item \textbf{Priority 1 (Rigid Resupply)}: Non-water maintenance materials ($L_{maint} = 3 \times 10^5$ MT/year) are subtracted from the total SES capacity $R_e$ first.
	\item \textbf{Priority 2 (Vital Consumables)}: Water is then allocated to the remaining SES capacity ($Cap_{avail}$).
\end{equation}
The \textbf{Capacity Overflow} occurs when $M_{net\_water} > Cap_{avail}$, forcing the system to utilize the costly Rocket System (RS):
\begin{equation}
	x_{water\_roc} = \max(0, M_{net\_water} - 237,000)
\end{equation}

\subsection{Performance Metrics: Marginal Cost and Infrastructure Duty Cycle}
To quantify the impact, we define two critical metrics:
\begin{itemize}
	\item \textbf{Marginal Cost ($\Delta C$)}: The step-function increase in annual budget when water resupply spills over from SES to RS rates ($C_{rocket} \gg C_{ele}$).
	\item \textbf{Occupancy Duty Cycle ($\Delta T$)}: The percentage of the SES annual schedule "locked" by water transport: $\Delta T = (x_{water\_ele} / R_e) \times 100\%$.
\end{itemize}

\subsection{Strategic Scenarios: Identifying the "Logistics Collapse" Tipping Point}
Our simulations reveal a critical \textbf{Tipping Point} based on recycling technology:
\begin{itemize}
	\item \textbf{The Efficiency Threshold}: At high efficiency ($\eta = 98\%$), the water load is a "minor pulse" ($Duty\_Cycle \approx 0.6\%$). However, as $\eta$ drops below a critical threshold, the $M_{net\_water}$ exceeds $Cap_{avail}$, triggering a \textbf{Logistics Collapse}.
	\item \textbf{The ISRU Imperative}: When $\Delta C$ exceeds the projected Colonial GDP, Earth-based supply becomes economically terminal. This mathematically justifies the necessity of \textbf{Lunar In-Situ Resource Utilization (ISRU)} not as an option, but as a prerequisite for survival.
\end{itemize}
	\section{Sensitivity Analysis}
	\lipsum[6]
	\section{Expansion of the Model}
	\section{Model Evaluation}
	\subsection{Strengths}
	\lipsum[4]
	\subsection{weaknesses}
	\begin{itemize}
		\item \textbf{Applies widely}\\
		This  system can be used for many types of airplanes, and it also
		solves the interference during  the procedure of the boarding
		airplane,as described above we can get to the  optimization
		boarding time.We also know that all the service is automate.
		\item \textbf{Improve the quality of the airport service}\\
		Balancing the cost of the cost and the benefit, it will bring in
		more convenient  for airport and passengers.It also saves many
		human resources for the airline.
	\end{itemize}
	
	AI cite use \AIcite{AI1,AI2,AI3}
	
	%\section{References}  % article类的section是最高级，自动进目录
	\cleardoublepage 
	\addcontentsline{toc}{section}{References}
	\begin{thebibliography}{99}
		\bibitem{gerald20xxwater}
		Gerald, J. (20XX). Water demand analysis for extraterrestrial habitats.
		Journal of Space Resource Utilization, XX(X), 1-15.
		
		\bibitem{wil20xxiss}
		Wil, X. (20XX). Water recycling efficiency of the International Space Station
		post-BPA upgrade. Acta Astronautica, XX(X), 89-102.
		
		\bibitem{isec2019mars}
		ISEC. (2019). Fast Transit to Mars: Apex Anchor Interplanetary Windows.
		ISEC Slide Deck.
		
		\bibitem{ipcc2014ar5}
		IPCC. (2014). Climate Change 2014: Mitigation of Climate Change.
		Contribution of Working Group III to the Fifth Assessment Report of the
		Intergovernmental Panel on Climate Change (AR5 WGIII). Cambridge University
		Press, Cambridge, UK and New York, NY, USA.
	\end{thebibliography}
	\begin{appendices}
		
		\section{First appendix}
		
		In addition, your report must include a letter to the Chief Financial Officer (CFO) of the Goodgrant Foundation, Mr. Alpha Chiang, that describes the optimal investment strategy, your modeling approach and major results, and a brief discussion of your proposed concept of a return-on-investment (ROI). This letter should be no more than two pages in length.
		
		\begin{letter}{Dear, Mr. Alpha Chiang}
			
			\lipsum[1-2]
			
			\vspace{\parskip}
			
			Sincerely yours,
			
			Your friends
			
		\end{letter}
		Here are simulation programmes we used in our model as follow.\\
		
		\textbf{\textcolor[rgb]{0.98,0.00,0.00}{Input matlab source:}}
		
		\section{Second appendix}
		
		some more text \textcolor[rgb]{0.98,0.00,0.00}{\textbf{Input C++ source:}}
		
	\end{appendices}
	\AImatter
	\begin{ReportAiUse}{9}
		\bibitem{AI1}
		Bing AI\\
		Query1: write a sort algorithm\\
		Output: Sure, I can help you with that. Here is a simple implementation of the selection sort algorithm in Python:
		\begin{lstlisting}[language=python]
			def selection_sort(arr):
			for i in range(len(arr)):
			min_idx = i
			for j in range(i + 1, len(arr)):
			if arr[min_idx] > arr[j]:
			min_idx = j
			arr[i], arr[min_idx] = arr[min_idx], arr[i]
			return arr
		\end{lstlisting}
		This algorithm sorts an array by repeatedly finding the minimum element from the unsorted part of the array and putting it at the beginning. \href{https://www.geeksforgeeks.org/sorting-algorithms/}{The time complexity of this algorithm is O($n^2$)}.
		
		There are many other sorting algorithms available such as bubble sort, insertion sort, merge sort, quick sort, heap sort, and more. Each of these algorithms has its own advantages and disadvantages. You can learn more about them on GeeksforGeeks and Programiz.
		
		I hope this helps!
		
		\bibitem{AI2}
		OpenAl Enie (Nov 5, 2023 version, Ernie 4.0)\\
		Query1: <insert the exact wording of any subsequent input into the Al tool>\\
		Output: <insert the complete output from the second query>
		
		\bibitem{AI3}
		Github CoPilot (Feb 3, 2024 version)\\
		Query1: <insert the exact wording you input into the Al tool>\\
		Output: <insert the complete output fron the Al tool>
		
		\bibitem{AI4}
		Google Bard (Feb 2,2024 version)\\
		Query: <insert the exact wording of your query>\\
		Output: <irsert the complete output from the Al tool>
		
		\bibitem{AI5}
		Bing AI\\
		Query: write a FFT algorithm\\
		Output: Sure, here is a Python implementation of the Fast Fourier Transform (FFT) algorithm:
		\begin{lstlisting}[language=python]
			import numpy as np
			
			def fft(x):
			n = len(x)
			if n == 1:
			return x
			even = fft(x[0::2])
			odd = fft(x[1::2])
			T = [np.exp(-2j * np.pi * k / n) * odd[k] for k in range(n // 2)]
			return [even[k] + T[k] for k in range(n // 2)] + [even[k] - T[k] for k in range(n // 2)]
		\end{lstlisting}
		The FFT is a widely used algorithm for computing the \textbf{Discrete Fourier Transform (DFT)} of a sequence. It reduces the number of computations needed from O(n$^2$) to O(n log n) 12. The basic idea of the FFT is to apply divide and conquer. We divide the coefficient vector of the polynomial into two vectors, recursively compute the DFT for each of them, and combine the results to compute the DFT of the complete polynomial 3.
		
		I hope this helps!
		
	\end{ReportAiUse}
	
\end{document}
%% 
%% This work consists of these files mcmthesis.dtx,
%%                                   figures/ and
%%                                   code/,
%% and the derived files             mcmthesis.cls,
%%                                   mcmthesis-demo.tex,
%%                                   README,
%%                                   LICENSE,
%%                                   mcmthesis.pdf and
%%                                   mcmthesis-demo.pdf.
%%
%% End of file `mcmthesis-demo.tex'.
